\documentclass[11pt,a4paper]{article}
\usepackage[hyperref]{acl2020}
\usepackage{times}
\usepackage{latexsym}
\usepackage{url}
\usepackage{booktabs}

\aclfinalcopy

\title{MuslimNews: A Dual-Track RAG and Instruction-Tuning Corpus of 173 International News Articles on the Global Muslim World}

\author{Salaamalykum Open Data Hub \\
  \texttt{https://salaamalykum.com} \\
  \texttt{https://github.com/salaamalykum/muslimnews}}

\date{2026-10-30}

\begin{document}
\maketitle
\begin{abstract}
We present \textbf{MuslimNews} (v1.0.0), a curated, cryptographically verified corpus of 173 Chinese-language international news articles and 1092 semantic chunks documenting global Muslim communities, public policy, religious heritage protection, and cross-cultural dialogues. Extracted and normalized from the WeCenter-powered Salaamalykum knowledge platform, MuslimNews implements a dual-track publishing architecture: human-readable Markdown with Schema.org JSON-LD on GitHub Pages, and page-indexed Apache Parquet alongside Croissant 1.0 metadata on Hugging Face.
\end{abstract}

\section{Introduction}
Multilingual retrieval-augmented generation (RAG) systems frequently suffer from factual blind spots when answering queries about contemporary events in the Islamic world phrased in Chinese. Existing news corpora lack domain-specific entity preservation and verifiable provenance. MuslimNews addresses this gap with 173 articles carrying SHA-256 content hashes, canonical cross-links, and multi-directional schema alignment with \texttt{quran\_rag\_dataset}.

\section{Corpus Architecture and Quality Gates}
Every document passes a seven-stage quality gate:
\begin{enumerate}
  \item Strict 8+ field YAML frontmatter (\texttt{original\_url}, \texttt{canonical\_url}, \texttt{author}, \texttt{pub\_date}, \texttt{updated\_at}, \texttt{language}, \texttt{topics}, \texttt{content\_hash}).
  \item Triple timestamp synchronization across site \texttt{lastmod}, RSS/Atom feeds, and release tags (2026-10-30T07:50:00Z).
  \item Dual-track separation between human-readable GFM (\texttt{content/}) and machine-trainable Parquet/JSONL views (\texttt{metadata/} and Hugging Face).
\end{enumerate}

\end{document}
